\section{Types}
\label{sec:types}

Types are a second static layer over the well-formed programs of \secref{well-formedness}. Every parameter,
return and assignment carries an annotation, so no type is inferred across a definition; within an
expression, types flow in the two directions of a bidirectional system. \emph{Synthesis} $\Gamma \vdash e
\synth \tau$ reads a type out of an expression, and \emph{checking} $\Gamma \vdash e \checks \tau$ tests an
expression against a type given by its surroundings (\figrefTwo{typing-synth}{typing-check}). A context now maps a variable to
its type, and definite assignment is the presence of a type: $\Gamma(x) = \tau$ says that $x$ is definitely
assigned and holds values of type $\tau$.

Variables, literals, module attributes, calls, constructor calls, operators and subscripts synthesise. A
literal has the type $\tyOf(\ell)$ of its form:
\begin{defalign}
\tyOf(n) & = \tyInt \quad \text{($n$ an integer literal)} \\
\tyOf(n) & = \tyFloat \quad \text{($n$ a float literal)} \\
\tyOf(w) & = \tyStr \\
\tyOf(\pyinline{True}) = \tyOf(\pyinline{False}) & = \tyBool \\
\tyOf(\pyinline{None}) & = \tyNone
\end{defalign}

\noindent The class entry $C$ is read as a map from field names to their declared types, so a constructor
call checks each argument against the type of the field it supplies and synthesises the class (\ruleName{syn-constr}).
Each operator has a signature $\odot : (\tau, \tau') \to \tau''$ relating its operand types to its result
type. Subscripting has a rule per container, and a tuple is subscripted by a literal index, since the
component types differ.

Displays, comprehensions and lambdas check. Their element types are not determined by the expression itself:
$\exList{}$ is a list of any element type, and a lambda has no annotations to declare its parameter types
(Python's grammar admits none), so both take their type from the surroundings. A tuple display and a
conditional expression have rules in both modes, synthesising when their parts do. One rule changes mode:
an expression that synthesises $\sigma$ checks against any supertype of $\sigma$ (\ruleName{chk-synth}).

Subtyping (\figref{subtyping}) is reflexivity, the union rules, and the ancestors of a class. Lists,
dictionaries, tuples and callables are related only by reflexivity: \pyinline{list} and \pyinline{dict} are
invariant in mypy, and \PurePy admits no program mypy rejects. A \pyinline{bool} is not an \pyinline{int},
though Python treats it as one.

\begin{figure}
\begin{subfigure}{\textwidth}
\flushleft \shadebox{$\Gamma \vdash e \synth \tau$}
\begin{mathpar}
\small
\inferrule*[lab={\ruleName{syn-var}}]
{
  \Gamma(x) = \tau
}
{
  \Gamma \vdash x \synth \tau
}
\and
\inferrule*[lab={\ruleName{syn-literal}}]
{
  \strut
}
{
  \Gamma \vdash \ell \synth \tyOf(\ell)
}
\and
\inferrule*[lab={\ruleName{syn-attr}}]
{
  \Gamma \vdash \exAttr{e}{x} : \tau
}
{
  \Gamma \vdash \exAttr{e}{x} \synth \tau
}
\and
\inferrule*[lab={\ruleName{syn-call}}]
{
  \Gamma \vdash e \synth \tyCallable{\tau_1, \ldots, \tau_n}{\tau} \\
  \Gamma \vdash e_i \checks \tau_i \quad (\forall i)
}
{
  \Gamma \vdash \exCall{e}{e_1, \ldots, e_n} \synth \tau
}
\and
\inferrule*[lab={\ruleName{syn-constr}}]
{
  \Gamma \vdash q : C \\
  \Gamma \vdash \fieldMap(C, \seq{e})(x) \checks C(x) \quad (\forall x \in \fields(C))
}
{
  \Gamma \vdash \exConstr{q}{\seq{e}} \synth q
}
\and
\inferrule*[lab={\ruleName{syn-binop}}]
{
  \odot : (\tau, \tau') \to \tau'' \\
  \Gamma \vdash e \checks \tau \\
  \Gamma \vdash e' \checks \tau'
}
{
  \Gamma \vdash \exBinOp{e}{\odot}{e'} \synth \tau''
}
\and
\inferrule*[lab={\ruleName{syn-unop}}]
{
  \ominus : \tau \to \tau' \\
  \Gamma \vdash e \checks \tau
}
{
  \Gamma \vdash \exUnOp{\ominus}{e} \synth \tau'
}
\and
\inferrule*[lab={\ruleName{syn-and}}]
{
  \Gamma \vdash e \checks \tyBool \\
  \Gamma \vdash e' \checks \tyBool
}
{
  \Gamma \vdash \exAnd{e}{e'} \synth \tyBool
}
\and
\inferrule*[lab={\ruleName{syn-or}}]
{
  \Gamma \vdash e \checks \tyBool \\
  \Gamma \vdash e' \checks \tyBool
}
{
  \Gamma \vdash \exOr{e}{e'} \synth \tyBool
}
\and
\inferrule*[lab={\ruleName{syn-subscript-list}}]
{
  \Gamma \vdash e \synth \tyList{\tau} \\
  \Gamma \vdash e' \checks \tyInt
}
{
  \Gamma \vdash \exSubscript{e}{e'} \synth \tau
}
\and
\inferrule*[lab={\ruleName{syn-subscript-str}}]
{
  \Gamma \vdash e \synth \tyStr \\
  \Gamma \vdash e' \checks \tyInt
}
{
  \Gamma \vdash \exSubscript{e}{e'} \synth \tyStr
}
\and
\inferrule*[lab={\ruleName{syn-subscript-dict}}]
{
  \Gamma \vdash e \synth \tyDict{\tau} \\
  \Gamma \vdash e' \checks \tyStr
}
{
  \Gamma \vdash \exSubscript{e}{e'} \synth \tau
}
\and
\inferrule*[lab={\ruleName{syn-subscript-tuple}}]
{
  \Gamma \vdash e \synth \tyTuple{\tau_1, \ldots, \tau_{n+1}} \\
  0 \leq i \leq n
}
{
  \Gamma \vdash \exSubscript{e}{i} \synth \tau_{i+1}
}
\and
\inferrule*[lab={\ruleName{syn-tuple}}]
{
  \Gamma \vdash e_i \synth \tau_i \quad (\forall i)
}
{
  \Gamma \vdash \exTuple{e_1, \ldots, e_n} \synth \tyTuple{\tau_1, \ldots, \tau_n}
}
\and
\inferrule*[lab={\ruleName{syn-cond}}]
{
  \Gamma \vdash e_1 \checks \tyBool \\
  \Gamma \vdash e \synth \sigma \\
  \Gamma \vdash e_2 \synth \tau
}
{
  \Gamma \vdash \exCondExpr{e_1}{e}{e_2} \synth \tyUnion{\sigma}{\tau}
}
\end{mathpar}
\subcaption{Expression $e$ synthesises type $\tau$ under $\Gamma$}
\end{subfigure}
\caption{Expression typing: synthesis}
\label{fig:typing-synth}
\end{figure}

\begin{figure}
\begin{subfigure}{\textwidth}
\flushleft \shadebox{$\Gamma \vdash e \checks \tau$}
\begin{mathpar}
\small
\inferrule*[lab={\ruleName{chk-synth}}]
{
  \Gamma \vdash e \synth \sigma \\
  \sigma \subty \tau
}
{
  \Gamma \vdash e \checks \tau
}
\and
\inferrule*[lab={\ruleName{chk-lambda}}]
{
  \Gamma \override \set{\bind{x_1}{\tau_1}, \ldots, \bind{x_n}{\tau_n}} \vdash e \checks \tau
}
{
  \Gamma \vdash \exLambda{x_1, \ldots, x_n}{e} \checks \tyCallable{\tau_1, \ldots, \tau_n}{\tau}
}
\and
\inferrule*[lab={\ruleName{chk-list}}]
{
  \Gamma \vdash e_i \checks \tau \quad (\forall i)
}
{
  \Gamma \vdash \exList{e_1, \ldots, e_n} \checks \tyList{\tau}
}
\and
\inferrule*[lab={\ruleName{chk-tuple}}]
{
  \Gamma \vdash e_i \checks \tau_i \quad (\forall i)
}
{
  \Gamma \vdash \exTuple{e_1, \ldots, e_n} \checks \tyTuple{\tau_1, \ldots, \tau_n}
}
\and
\inferrule*[lab={\ruleName{chk-dict}}]
{
  \Gamma \vdash e_i \checks \tyStr \quad (\forall i) \\
  \Gamma \vdash e_i' \checks \tau \quad (\forall i)
}
{
  \Gamma \vdash \exDict{\bind{e_1}{e_1'}, \ldots, \bind{e_n}{e_n'}} \checks \tyDict{\tau}
}
\and
\inferrule*[lab={\ruleName{chk-list-comp}}]
{
  \Gamma \vdash \seq{g} : \Delta \\
  \Gamma \override \Delta \vdash e \checks \tau
}
{
  \Gamma \vdash \exListComp{e}{g_1\;\ldots\;g_n} \checks \tyList{\tau}
}
\and
\inferrule*[lab={\ruleName{chk-dict-comp}}]
{
  \Gamma \vdash \seq{g} : \Delta \\
  \Gamma \override \Delta \vdash e \checks \tyStr \\
  \Gamma \override \Delta \vdash e' \checks \tau
}
{
  \Gamma \vdash \exDictComp{e}{e'}{g_1\;\ldots\;g_n} \checks \tyDict{\tau}
}
\and
\inferrule*[lab={\ruleName{chk-cond}}]
{
  \Gamma \vdash e_1 \checks \tyBool \\
  \Gamma \vdash e \checks \tau \\
  \Gamma \vdash e_2 \checks \tau
}
{
  \Gamma \vdash \exCondExpr{e_1}{e}{e_2} \checks \tau
}
\end{mathpar}
\subcaption{Expression $e$ checks against type $\tau$ under $\Gamma$}
\end{subfigure}

\vspace{2mm}

\begin{subfigure}{\textwidth}
\flushleft \shadebox{$\Gamma \vdash \seq{g} : \Delta$}
\begin{mathpar}
\small
\inferrule*[lab={\ruleName{ty-quals-nil}}]
{
  \strut
}
{
  \Gamma \vdash \varepsilon : \varnothing
}
\and
\inferrule*[lab={\ruleName{ty-qual-guard}}]
{
  \Gamma \vdash e \checks \tyBool \\
  \Gamma \vdash \seq{g} : \Delta
}
{
  \Gamma \vdash \qualGuard{e} \cons \seq{g} : \Delta
}
\and
\inferrule*[lab={\ruleName{ty-qual-generator}}]
{
  \Gamma \vdash e \synth \tau \\
  \elemTy(\tau) = \tau' \\
  \Gamma \override \set{\bind{x}{\tau'}} \vdash \seq{g} : \Delta
}
{
  \Gamma \vdash \qualGenerator{x}{e} \cons \seq{g} : \set{\bind{x}{\tau'}} \override \Delta
}
\end{mathpar}
\subcaption{Comprehension qualifiers $\seq{g}$ bind $\Delta$ under $\Gamma$}
\end{subfigure}
\caption{Expression typing: checking}
\label{fig:typing-check}
\end{figure}


\begin{figure}
\flushleft \shadebox{$\Gamma \vdash \sigma \subty \tau$}
\begin{mathpar}
\small
\inferrule*[lab={\ruleName{subty-refl}}]
{
  \strut
}
{
  \Gamma \vdash \tau \subty \tau
}
\and
\inferrule*[lab={\ruleName{subty-class}}]
{
  \Gamma(c') \in \ancestors(\Gamma(c))
}
{
  \Gamma \vdash q.c \subty q.c'
}
\and
\inferrule*[lab={\ruleName{subty-never}}]
{
  \strut
}
{
  \Gamma \vdash \tyNever \subty \tau
}
\and
\inferrule*[lab={\ruleName{subty-object}}]
{
  \strut
}
{
  \Gamma \vdash \tau \subty \tyObject
}
\and
\inferrule*[lab={\ruleName{subty-int-float}}]
{
  \strut
}
{
  \Gamma \vdash \tyInt \subty \tyFloat
}
\and
\inferrule*[lab={\ruleName{subty-literal}}]
{
  \Gamma \vdash \widen(\ell) \subty \tau
}
{
  \Gamma \vdash \tyLiteral{\ell} \subty \tau
}
\and
\inferrule*[lab={\ruleName{subty-union-left}}]
{
  \Gamma \vdash \sigma \subty \tau
}
{
  \Gamma \vdash \sigma \subty \tyUnion{\tau}{\tau'}
}
\and
\inferrule*[lab={\ruleName{subty-union-right}}]
{
  \Gamma \vdash \sigma \subty \tau'
}
{
  \Gamma \vdash \sigma \subty \tyUnion{\tau}{\tau'}
}
\and
\inferrule*[lab={\ruleName{subty-union}}]
{
  \Gamma \vdash \tau \subty \sigma \\
  \Gamma \vdash \tau' \subty \sigma
}
{
  \Gamma \vdash \tyUnion{\tau}{\tau'} \subty \sigma
}
\end{mathpar}
\caption{Subtyping}
\label{fig:subtyping}
\end{figure}


A generator draws its bindings from the elements of a value, and $\elemTy$ gives the type of those elements,
mirroring $\iterOp$ (\secref{auxiliary-values}):
\begin{defalign}
\elemTy(\tyList{\tau}) & = \tau \\
\elemTy(\tyStr) & = \tyStr \\
\elemTy(\tyDict{\tau}) & = \tyStr \\
\elemTy(\tyTuple{\tau, \ldots, \tau}) & = \tau
\end{defalign}

\noindent A tuple of mixed component types has no element type, so it cannot be iterated.
